\fancyhead[L]{\small PASQUICH ON SPHERICAL ASTRONOMY.}
\SourceLeaf{540}
\begin{center}\small Göttingische gelehrte Anzeigen. No. 144. pp. 1433..1436. September 9, 1811.\\[3pt]\rule{.72\linewidth}{.3pt}\end{center}

{\rightskip=0pt plus 1em \emph{Compendium of the Elements of Spherical and Computational Astronomy by} \textsc{Joannes Pasquich}, director of the astronomical observatory of the Royal Hungarian University. \emph{Part One. {\addfontfeatures{LetterSpace=10}Theoretical} Elements of Spherical and Computational Astronomy. Part Two, {\addfontfeatures{LetterSpace=10}Practical} Elements of Spherical and Computational Astronomy.} 1811. 160 and 166 pages in quarto. Vienna, Schaumann and Co.\par}

In the preface the author explains that the book grew out of lectures on astronomy given at the University of Buda. He wrote it for people just starting out, not for experienced astronomers. The book aims less to give beginners a complete guide to all the various procedures of astronomy than to draw them to the study of a science that has so far been too neglected there and encourage them to pursue it, so that they can go on to receive more advanced instruction. This, he says, should excuse the otherwise striking scantiness of the book. It is to the author's credit that he seeks to achieve this aim by starting with a thorough grounding in the fundamentals instead of skimming off entertaining results (which does nothing to further serious study). The first, theoretical part covers the following topics in five sections. The \emph{first} section, entitled “Basic Concept and Fundamental Laws of the Daily Motion of the Heavenly Bodies and the Annual Motion of the Sun,” is divided into four chapters. These introduce the basic concepts of the sphere and the circles drawn on it, along with the trigonometric formulas for determining the positions of the heavenly bodies relative to those circles. The \emph{second} section, also divided into four chapters, deals with the different kinds of time and the related use of right ascensions, as well as differences of meridian. Given the thoroughness of the exposition everywhere else, one might wish that a point or two here had been taken a little further. For example, the notion of mean solar time as a \emph{duration} is quite clear, but the notion of mean solar time as an \emph{instant} is still uncertain. One cannot see from what point the mean and true motions of the Sun are to be assumed to start out together. The \emph{third} section, in three chapters, deals with the relations between the motion of the Sun in the ecliptic and its right ascensions, declinations and position angles. It also deals with the altitudes, hour angles and azimuths of the heavenly bodies. The \emph{fourth} section discusses the shape and size of the Earth, parallax and astronomical refraction in four chapters. In the five chapters of the \emph{fifth} section the author teaches the rudiments of the elliptical motion of the planets and the phenomena that depend on it, and the effects of precession, nutation and aberration. This concludes the theoretical part. This account of the contents shows that the author aimed more at providing a collection of useful formulas for calculation than at arranging the exposition so that readers could gain an instructive and satisfying understanding, without further oral help, of how the truths of astronomy were discovered or could be discovered. Much therefore has to be left to such help.

The \emph{second part} of the work can by no means be regarded as a complete handbook of computational astronomy. It confines itself to the practical rules that an observing astronomer is likely to need most often. It contains three sections. The \emph{first} explains in practical terms how astronomical tables in general, and solar tables in particular, are arranged, how to use the astronomical\SourceLeaf{541} ephemerides and catalogues of fixed stars, and how to calculate nutation and aberration. The \emph{second} section teaches how to compare the different kinds of astronomical time and how to determine time from observed culminations and corresponding altitudes. The \emph{third} section shows chiefly how to use observed circummeridian altitudes to determine latitude and describes how to derive the positions of heavenly bodies from observations.

The work is accompanied by an Appendix to Joannes Pasquich's Compendium of the Practical Elements of Spherical and Computational Astronomy, Containing Auxiliary Tables. Vienna, Schaumburg and Co. 1811. 42 pages in quarto. The tables comprise \textsc{Delambre}'s solar tables in abridged form, the general table for the correction for noon in its usual arrangement, an auxiliary table for reducing zenith distances measured off the meridian, another auxiliary table for reducing the zenith distances of the Pole Star, a table for determining precession in right ascension and declination, tables for refraction, aberration and nutation, and \textsc{Piazzi}'s latest catalogue of 121 fundamental stars. The reviewer would have liked an explanation to accompany some of the tables whose arrangement cannot be understood on its own (such as the tables for reducing zenith distances), since some readers might want to buy the tables without the Compendium.
